\section{Reduction of the twin conjecture to a single node}\label{sec:reduction}

This section is the spine of the article. We compress the entire twin branch into one
machine-verified chain and expose, at its far end, the single open node on which everything hangs
--- and the single \code{sorry}-tainted declaration \code{twin\_prime\_conjecture} that records it.
Throughout we insist on the discipline of the introduction: this is a \emph{reduction}, not a proof.
Every transit step below is \green{} machine-checked and distribution-free; all the difficulty is
quarantined in one \red{} statement, and we never pass the reduction off as a closure of the
conjecture.

\subsection{The bridge to infinity}

We work with the centre form of a twin pair. Every twin pair beyond $(3,5)$ has the shape
$(6m-1,6m+1)$, so the centre encoding loses nothing.

\begin{definition}[\code{IsTwinCenter}, \code{TwinLowers}]\label{def:istwincenter}
For $m\in\N$,
\[
  \mathrm{IsTwinCenter}(m)\;:\Longleftrightarrow\;(6m-1)\ \text{prime}\ \wedge\ (6m+1)\ \text{prime},
\]
and $\mathrm{TwinLowers}=\{\,p\mid p\ \text{and}\ p+2\ \text{prime}\,\}$. The twin conjecture is
exactly \code{TwinLowers.Infinite}.
\end{definition}

The bridge from block combinatorics to arithmetic is assembled from three elementary, non-circular
links. The first is pure cardinality.

\begin{theorem}[\code{survivor\_of\_not\_covered}]\label{thm:survivor_of_not_covered}
For finite $\Omega,\mathrm{bad}\subseteq\N$,
\begin{equation}\label{eq:survivor}
  |\mathrm{bad}|<|\Omega|\;\Longrightarrow\;\exists\,m\in\Omega,\ m\notin\mathrm{bad}.
\end{equation}
\end{theorem}

\emph{Proof idea.} The pigeonhole principle in its purest form: were $\Omega\subseteq\mathrm{bad}$,
monotonicity of cardinality (\code{Finset.card\_le\_card}) would give $|\Omega|\le|\mathrm{bad}|$,
against the premise. No analysis, no distribution of primes --- only \code{Finset} and the order on
$\N$. \green{}

The second link is the classical primality criterion, and its immediate corollary for a centre.

\begin{theorem}[\code{prime\_of\_no\_small\_prime\_factor}]\label{thm:prime_of_no_small_prime_factor}
Let $n\ge 2$ and suppose $\forall p\ \text{prime},\ p^2\le n\Rightarrow p\nmid n$. Then $n$ is prime.
\end{theorem}

\emph{Proof idea.} By contradiction through the least prime divisor $p=\mathrm{minFac}(n)$:
minimality gives $p\le n/p$ (\code{Nat.minFac\_le\_div}), whence $p^2\le n$, contradicting the
hypothesis. A composite number's least prime divisor never exceeds $\sqrt n$, so a sieve to the root
settles primality. \green{}

\begin{theorem}[\code{isTwinCenter\_of\_root\_sieve}]\label{thm:isTwinCenter_of_root_sieve}
If $6m\pm1\ge 2$ and no prime up to the root of either side divides that side, then
$\mathrm{IsTwinCenter}(m)$.
\end{theorem}

\emph{Proof idea.} Two applications of \cref{thm:prime_of_no_small_prime_factor}, one per side,
packaged into the conjunction \code{IsTwinCenter}. The direction \emph{survivor $\Rightarrow$ twin}
carries no hypothesis at all: it is elementary, and thus fixes an important boundary --- all the
substantive difficulty lies not here but in \emph{producing} a survivor sieved to the root. \green{}

The third link is order-theoretic.

\begin{theorem}[\code{infinite\_of\_unbounded\_centers}]\label{thm:infinite_of_unbounded_centers}
If for every $N$ there is a centre $m>N$ with $\mathrm{IsTwinCenter}(m)$, then
\code{TwinLowers.Infinite}.
\end{theorem}

\emph{Proof idea.} By \code{Set.infinite\_of\_not\_bddAbove}: the witness $6m-1\in\mathrm{TwinLowers}$
exceeds $N$ and its neighbour $6m+1$ is prime, so \code{TwinLowers} is unbounded above, hence
infinite. Monotonicity of $6m-1$ in $m$ turns ``arbitrarily high centre'' into ``arbitrarily large
lower twin''. \green{}

\subsection{Isolating the block core}

The three links glue into a local transport lemma and a global reduction. Fix, at scale $N$, a
finite \emph{carrier} of candidate centres above $N$ and a finite \emph{bad} set of centres that are
certainly not twin centres.

\begin{theorem}[\code{twin\_center\_of\_block}]\label{thm:twin_center_of_block}
Let $\mathit{carrier},\mathit{bad}\subseteq\N$ be finite with
\begin{equation}\label{eq:block}
  (\forall m\in\mathit{carrier},\ N<m),\quad |\mathit{bad}|<|\mathit{carrier}|,\quad
  (\forall m\in\mathit{carrier},\ m\notin\mathit{bad}\Rightarrow\mathrm{IsTwinCenter}(m)).
\end{equation}
Then $\exists m,\ N<m\ \wedge\ \mathrm{IsTwinCenter}(m)$.
\end{theorem}

\emph{Proof idea.} \cref{thm:survivor_of_not_covered} extracts a survivor $m\in\mathit{carrier}$ with
$m\notin\mathit{bad}$; the first conjunct places it above $N$, the third makes it a twin centre. Two
lines of \code{Finset} combinatorics, no appeal to prime density. \green{}

\begin{theorem}[\code{twin\_prime\_conjecture\_of\_blocks}]\label{thm:twin_prime_conjecture_of_blocks}
If for every $N$ there exist finite $\mathit{carrier},\mathit{bad}$ satisfying \cref{eq:block}, then
\code{TwinLowers.Infinite}.
\end{theorem}

\emph{Proof idea.} \cref{thm:infinite_of_unbounded_centers} reduces the goal to the unboundedness of
twin centres; for each $N$ the block premise plus \cref{thm:twin_center_of_block} supplies a twin
centre above $N$. \green{} The implication is non-circular: it nowhere uses the conjecture, nor hides
it inside the premise. What remains open is exactly the premise --- a strict cardinality inequality
plus a survival criterion, that is, a purely combinatorial problem about finite sets. Node~(3),
survivor $\Rightarrow$ twin, is already closed by \cref{thm:isTwinCenter_of_root_sieve}; the
\emph{counting core} (that the carrier outweighs the bad set) is what remains.

\subsection{The first named node: the real four-corner}

The counting core first crystallises as the strict \emph{real four-corner} on rank counts. Writing
$R_{ij}$ for the rank sets, the hypothesis $H$ (\code{FourCornerInput}) asks, at every scale $N$:
$0<|R_{00}|$, the strict four-corner $|R_{00}|\,|R_{33}|<|R_{03}|\,|R_{30}|$, the light side-corner
$|R_{03}|\,|R_{30}|\le|R_{00}|^2$, together with $|\mathit{carrier}|=|R_{00}|$,
$|\mathit{bad}|=|R_{33}|$ and the carrier/twin conditions of \cref{eq:block}. Via
\code{N33\_lt\_N00\_of\_four\_corner} these yield $|R_{33}|<|R_{00}|$, i.e. exactly
$|\mathit{bad}|<|\mathit{carrier}|$.

\begin{theorem}[\code{twin\_primes\_of\_four\_corner}]\label{thm:twin_primes_of_four_corner}
\green{} (a green conditional theorem: axiom-clean, conditional on the named input $H$)
\begin{equation}\label{eq:fourcorner}
  H\;\Longrightarrow\;\code{TwinLowers.Infinite}.
\end{equation}
\end{theorem}

\emph{Proof idea.} Unfold $H\,N$, derive $|\mathit{bad}|<|\mathit{carrier}|$ from the four-corner,
and apply \cref{thm:twin_prime_conjecture_of_blocks}. There is no \code{sorry}: the conclusion follows
from the explicit input $H$, and all nontrivial content is carried by $H$. Contraposition records the
same fact by hand.

\begin{theorem}[\code{twin\_finite\_contradiction}]\label{thm:twin_finite_contradiction}
If $\neg\,\code{TwinLowers.Infinite}$ and \code{FourCornerInput} $H$, then \code{False}.
\end{theorem}

\emph{Proof idea.} Literally $\mathrm{hfin}\,(\code{twin\_primes\_of\_four\_corner}\,H)$: finiteness is
a map from infinitude into the void, fed exactly the infinitude the four-corner manufactures.
\green{} Logically contraposition adds nothing, but it makes tangible what breaks when the thesis is
denied: a finite tail must be covered by bad centres with no survivor, which $H$ forbids.

\begin{remark}
Three independent routes --- a union bound over primes, EPMI infinite descent, and the four-corner
--- were probed, and all reduce to the same object: the strict real four-corner plus a lower bound on
the carrier. The margin $|R_{03}|\,|R_{30}|-|R_{00}|\,|R_{33}|$, normalised by scale, numerically
creeps toward zero: the model four-corner (\code{model\_four\_corner}, proved unconditionally) does
\emph{not} transfer to reality distribution-free. This is the parity problem, and it is quarantined
entirely inside the \red{} input $H$.
\end{remark}

\subsection{The payment ledger and the defect law}

Halting is not free: it has an algebraic price, provable with only \code{Nat}/\code{Int} and
\code{Finset}. A free passage to an ordered prime $p$ requires tax-freedom at every smaller small
prime, so the shift $a-\theta$ (with $\theta=\sigma\varepsilon$) must be divisible by the whole
primorial $P$ of those primes.

\begin{theorem}[\code{shifted\_primorial\_bound}]\label{thm:shifted_primorial_bound}
If $P\mid a-\theta$ and $a\ne\theta$, then $P\le|a-\theta|$.
\end{theorem}

\emph{Proof idea.} From $a\ne\theta$ we have $a-\theta\ne0$, so $|a-\theta|>0$ and
\code{Int.le\_of\_dvd} gives $P\le|a-\theta|$. \green{} This is the \emph{defect law}: once the
primorial outgrows the active divisor, there is no nontrivial tax-free passage
(\code{late\_boundary\_not\_free}). It maps the parity wall precisely but does not bypass it --- the
total capacity loss $\prod_{q\le A}\tfrac{q-2}{q-1}\sim 1/\ln A$ (Mertens) is a wall that no
distribution-free bookkeeping breaks. The single scalar input of this route is again the same
parity/sieve-remainder obstruction.

\subsection{SNOL: reduction to the rank-1 neighbour}

Rather than defeat the counting budget, we reduce the whole programme to one algebraic node that is
\emph{non-counting by construction}. Defects of a product-state are organised by rank, and rank
descent reduces every defect to rank~$1$, where one side of the wedge is the active prime $a$ itself,
so the opposite side is the shifted neighbour $a-2\varepsilon$. The rank-1 algebra
(\code{rank1\_opposite}, \code{neighbour\_saturation}, \code{neighbour\_corridor\_bound}) is proved
without a single distributional assumption. Feeding the same block bridge yields the conditional
theorem.

\begin{theorem}[\code{twin\_primes\_of\_SNOL}]\label{thm:twin_primes_of_SNOL}
If \code{SNOLInput} holds --- at every scale $N$ a carrier above $N$ whose bad centres (carriers of the
terminal shifted-neighbour obstruction $p\mid a-2\varepsilon$, $p\le A$) are strictly fewer, every
survivor a twin --- then \code{TwinLowers.Infinite}.
\end{theorem}

\emph{Proof idea.} The direct bridge \cref{thm:twin_prime_conjecture_of_blocks} with the input
supplied as SNOL rather than the four-corner; \code{finite\_contradicts\_SNOL} completes the
contraposition. \green{} conditional on the \red{} node \code{SNOLInput}. The neighbour audit shows
$62$--$78\%$ of active primes already have a small-prime-caught neighbour, and the fraction
\emph{grows} with $A$: SNOL cannot be proved or refuted by counting/CRT capacity in either direction.
It must therefore rest on the \emph{Euclidean pedigree} of $a$ --- that $a$ arrived from the descent
of a wedge centre.

\subsection{The catch is a step of descent: old-peel closes onto the engine}

The terminal $p\mid a-2\varepsilon$ is \emph{not} a terminal. It is literally the divisibility of the
opposite side of the wedge (\code{catch\_is\_opposite}), and it unfolds into an \emph{old-peel}: a
decomposition $6n-\varepsilon=p\,(6t+\delta)$ producing a strictly smaller centre $t$, with a rigid
sign law $\delta=-\pi\varepsilon$ (\code{old\_peel\_sign}).

\begin{theorem}[\code{old\_peel\_height\_drop}]\label{thm:old_peel_height_drop}
If $6n-\varepsilon=p\,(6t+\delta)$ with $p\ge5$, $\varepsilon,\delta\in\{\pm1\}$, $t\ge1$, $n\ge2$,
then $t<n$.
\end{theorem}

\emph{Proof idea.} From $6t+\delta=(6n-\varepsilon)/p\le(6n+1)/5$ and $n\ge2$ we get $t<n$
(\code{nlinarith}). \green{} Numerically the descent is sharper still --- $t<n/5$ on $3000/3000$ real
rank-1 catches. Every catch is a step downward; the flow cannot ride down forever.

\begin{theorem}[\code{no\_infinite\_old\_peel}]\label{thm:no_infinite_old_peel}
For any $z:\N\to\N$ with \code{StrictAnti}\,$z$, \code{False}.
\end{theorem}

\emph{Proof idea.} This is \code{no\_infinite\_engine\_descent} applied to the old-peel height chain:
an infinite strictly descending chain in $\N$ does not exist. \green{} The final SNOL node thus closes
\emph{not onto a new distribution theorem but onto the already-proven engine}. What remains is the
structural premise that a non-sink centre always regenerates a correct successor.

\subsection{NOPSL: the flow reaches a twin}

We abstract exactly the structure old-peel supplies.

\begin{definition}[\code{OldPeelLedger}]\label{def:oldpeelledger}
An old-peel ledger over states $\sigma$ is a tuple $(h,\mathrm{Sink},\mathrm{Step},\code{step\_drops},
\code{regenerate})$ with $h:\sigma\to\N$, $\mathrm{Sink},\mathrm{Step}$ predicates, the old-peel law
$\code{step\_drops}:\mathrm{Step}\,s\,s'\Rightarrow h(s')<h(s)$, and regeneration
$\code{regenerate}:\neg\,\mathrm{Sink}(s)\Rightarrow\exists s',\ \mathrm{Step}\,s\,s'$.
\end{definition}

The structure knows nothing of primes or divisibility. All the arithmetic is spent on the two law
fields; the reasoning that follows is pure ``height falls'' and ``the flow does not hang''.

\begin{theorem}[\code{twin\_primes\_of\_nopsl}]\label{thm:twin_primes_of_nopsl}
Given a family of ledgers $L(N)$ with starts and a centre map, and the condition
$\code{sink\_is\_twin}:\ L(N).\mathrm{Sink}(s)\Rightarrow N<\mathrm{center}(N,s)\wedge
\mathrm{IsTwinCenter}(\mathrm{center}(N,s))$, we have \code{TwinLowers.Infinite}.
\end{theorem}

\emph{Proof idea.} By \code{no\_old\_peel\_sink}: assuming no sink, \code{regenerate} plus choice builds
an infinite trajectory whose heights strictly decrease (\code{step\_drops}), contradicting
\cref{thm:no_infinite_old_peel}. A sink then exists at every scale; \code{sink\_is\_twin} makes it a
twin centre above $N$, and \cref{thm:infinite_of_unbounded_centers} concludes. \green{} The full
implication $\code{OldPeelLedger}\Rightarrow\code{TwinLowers.Infinite}$ is machine-checked on the pure
logic of strict descent. Of the two law fields, \code{step\_drops} is proved arithmetic; the single
remaining input is the structural field \code{regenerate} --- ``the flow has nowhere to go except
downward or into a twin''.

\subsection{Regeneration is elementary in three of four classes}

The word ``regenerates'' hides an unconditional dichotomy of any centre $t$.

\begin{theorem}[\code{regeneration\_dichotomy}]\label{thm:regeneration_dichotomy}
For every centre $t$ and threshold $A$, exactly one holds:
\[
  \mathrm{Twin}(t)\ \lor\
  \bigl(\neg\mathrm{OldFree}(A,t)\wedge\exists q\ \text{prime},\,5\le q\le A,\,q\mid 6t\pm1\bigr)\ \lor\
  \bigl(\mathrm{OldFree}(A,t)\wedge(\neg(6t-1)\ \text{prime}\ \lor\ \neg(6t+1)\ \text{prime})\bigr).
\]
\end{theorem}

\emph{Proof idea.} \code{by\_cases} on $\mathrm{OldFree}\,A\,t$: not-old-free yields an old-peel divisor
(\code{not\_oldfree\_gives\_peel}); old-free splits by primality, and a composite old-free side carries
a divisor $>A$ (\code{composite\_oldfree\_has\_big\_divisor}), an active-descent edge. \green{}
Complete and unconditional. In classes (2),(3) an outgoing descent edge is exhibited; in (1) a
legitimate twin sink. The nontrivial remainder is class (4), fan-in/Hall --- a statement about the
\emph{set} of pedigrees, not derivable from $\mathrm{OldFree}$.

The parity/density inputs that once looked distributional are discharged \emph{existentially}: a clean
start above any $N$ is written by the formula $m=(N+1)P_A$ (\code{carrier\_nonempty\_above}), and a
clean sink below $A^2$ is automatically a twin.

\begin{theorem}[\code{sink\_is\_twin}]\label{thm:sink_is_twin}
If both sides satisfy $2\le 6m\pm1<A^2$ and no prime $q\le A$ divides either side, then
$\mathrm{TwinCenterZ}\,m$.
\end{theorem}

\emph{Proof idea.} A number without prime divisors $\le A$ and below $A^2$ cannot be composite (its
least prime divisor would satisfy $p^2\le n<A^2$, so $p<A$). Two copies of
\code{oldfree\_below\_sq\_prime}, one per side. \green{} Cleanliness excludes small divisors, the
threshold $A^2$ excludes large ones; no room is left for a composite side, and no prime density is
invoked.

\subsection{The clean graph: closing the sink-is-clean gap}

A descent from a clean start may reach a small unclean twin ($18\mapsto(107,109)$) to which the
anchoring lemma does not apply. The repair is definitional: a sink is defined only inside the clean
graph, and an edge to an unclean target is a \emph{boundary exit}, dispatched to the defect ledger.

\begin{theorem}[\code{clean\_center\_outcome}]\label{thm:clean_center_outcome}
For every clean centre $m$, exactly one holds:
\[
  \mathrm{TwinCenterZ}\,m\ \lor\ (\exists n,\ \mathrm{CleanActiveEdge}\,A\,m\,n)\ \lor\
  (\exists n,\ \mathrm{BoundaryExit}\,A\,m\,n)\ \lor\ \mathrm{CleanSink}\,A\,m.
\]
\end{theorem}

\emph{Proof idea.} Nested case analysis on twin-ness and on the existence of an active edge, closed by
\code{active\_edge\_clean\_or\_boundary} (excluded middle on $\mathrm{Clean}\,A\,n$). \green{} The
classification correctly excludes the unclean small twin from the sink case: $18$ settles in the
boundary branch, not the sink. But the graph is far from closed --- $59\%$ of clean centres have all
their active edges in the boundary. The whole remainder concentrates in one \red{} structural input:
a boundary state must regenerate.

\subsection{The global node and the pigeonhole skeleton}

Pointwise boundary regeneration is \emph{false} ($18\mapsto(107,109)$ has no successor yet is no
correct sink). What is provable pointwise is a dichotomy of outcome
(\code{boundary\_exit\_decomposes}): an impure target either continues the old-peel descent or settles
into a finite old absorber below the threshold $M_0$. The load then moves to the whole set of fresh
starts. If there are no new twins above $M_0$, infinitely many clean starts must settle into finitely
many old absorbers --- a fan-in the numbers record as up to $570\to1$.

\begin{theorem}[\code{global\_absorber\_forces\_engine}]\label{thm:global_absorber_forces_engine}
Let $S$ be infinite, $\mathrm{Codom}$ finite, and $\mathrm{key}:\alpha\to\mathrm{Codom}$ a start's
passport. Given the pump
\begin{equation}\label{eq:pump}
  \forall\gamma_1\gamma_2,\ \gamma_1\ne\gamma_2\to\mathrm{key}\,\gamma_1=\mathrm{key}\,\gamma_2\to\code{Engine},
\end{equation}
then \code{Engine}.
\end{theorem}

\emph{Proof idea.} An honest pigeonhole: without an engine, \code{pump} is forbidden, so \code{key} is
injective on $S$, so an infinite set maps injectively into a finite type
(\code{Set.Finite.of\_finite\_image}) --- impossible. The collision $\gamma_1\ne\gamma_2$ with equal
passport is the \green{} proven part; all the remaining load is in the \red{} hypothesis
\cref{eq:pump}, the Hall node ``two distinct genealogies with one passport $\Rightarrow$ engine''.
Assembled with the proven impossibility of the engine, \code{no\_global\_absorption\_under\_epmi} then
forces a new twin above $M_0$ --- conditional on the node.

\begin{remark}
Between chapters the boundary node was attacked from many angles --- weighted debt energy, inverse-limit
towers, jump/cut barriers, paid dynamics, a closed-universe hygiene, and a $\mathrm{P}$-versus-$\mathrm{NP}$
reading. Every wrapper was machine-audited; each either turned circular (the reduction became
equivalent to the goal), vacuous (its premise unreachable), or misoriented (all arithmetic dynamics run
downward while the promotion runs up). No wrapper lowered the content below the node, and two internal
audits caught genuine vacuity episodes before they reached the showcase. The red line held: the only
non-standard axiom is \axiom{}, and the only \code{sorry} is \code{twin\_prime\_conjecture}.
\end{remark}

\subsection{Rigid closure: descent without a cycle}

The old argument ``no sinks $\Rightarrow$ directed cycle $\Rightarrow$ engine $\Rightarrow\bot$'' is
logically redundant. Over a \code{HeightGraph} (a height, a twin predicate, a step with
$\code{step\_drops}:\mathrm{Step}\,s\,t\Rightarrow h(t)<h(s)$), strict fall of the height rules out an
infinite chain directly.

\begin{theorem}[\code{reaches\_twin}]\label{thm:reaches_twin}
If $G$ regenerates (every non-twin state has a descending successor), then
$\forall s,\ \exists t,\ \mathrm{Twin}\,t$.
\end{theorem}

\emph{Proof idea.} Strong induction on the height of the start: a non-twin state steps down to a
strictly smaller height, and one cannot descend forever in $\N$. \green{} The cycle branch and its
rigid signature vanish as separate inputs; \code{no\_cycle} keeps the impossibility of a cycle only as
a repackaging of well-foundedness. It remains to build the abstract edge by arithmetic.

\begin{theorem}[\code{cofactor\_is\_center}]\label{thm:cofactor_is_center}
For $t\ge1$, $\eta\in\{\pm1\}$, a prime $q\ge5$ coprime to $6$ with $q\mid 6t+\eta$, there exist
$t',\eta'$ with $\eta'\in\{\pm1\}$, $6t'+\eta'=(6t+\eta)/q$ and $0\le t'<t$.
\end{theorem}

\emph{Proof idea.} The cofactor $c=(6t+\eta)/q$ is forced coprime to $6$ (both $6t+\eta$ and $q$ are,
so $c\bmod6\in\{1,5\}$), positive, and $<t$ since $5c\le 6t+1$. \green{} Checked at $100\%$ on
$307010$ cases. Divisibility thus converts automatically into a valid smaller centre --- the edge is
\emph{built}, not postulated. The single remaining input is the completeness of the dichotomy on the
carrier scale (\code{regenerates\_needs\_target\_center}): that fan-in/Hall introduces no
unclassifiable terminal.

\subsection{The product-core machine: rank descent \texorpdfstring{$4\to1$}{4->1}}

The separating scale $6X_A+1<P_A$ (the primorial outruns the carrier bound, attainable with
exponential margin already at $A\approx50$) makes the coarse passport $a\bmod P_A$ injective on the
legal layer. Carried into rank induction, this closes \code{no\_productHall} on a tail-free
\code{RankNode} $r$ --- a sign plus $\mathrm{factors}:\mathrm{Fin}\,r\to\N$ --- where extensionality is
a genuine theorem (\code{no\_mismatch\_core\_eq}), not a hypothesis.

\begin{theorem}[\code{no\_productHall}]\label{thm:no_productHall}
If at role $k$ both factors are $<P_A$ and congruent mod $P_A$, yet distinct, then \code{False}.
\end{theorem}

\emph{Proof idea.} For $a<P_A$ we have $\code{ZMod.val}(a\bmod P_A)=a$, so equal residues force equal
factors, against the mismatch. \green{} Removing the tail from the induction object removes the
$570\to1$ fan-in, and with it all three audit defects (extensionality, the residual bound
\code{ambientLegal\_delete}, finiteness of the signature at fixed $A$) at once.

\begin{theorem}[\code{core\_step\_proved}]\label{thm:core_step_proved}
Under the separating scale, $\forall r,\ 2\le r\Rightarrow\code{CoreCollision}\,X_A\,P_A\,r\Rightarrow
\code{CoreCollision}\,X_A\,P_A\,(r-1)$.
\end{theorem}

\emph{Proof idea.} From equal signatures and distinct nodes there is a mismatched role $k$; delete it.
Either the residuals differ (a lower-rank collision) or they agree, concentrating the mismatch at $k$
--- a ProductHall, impossible by \cref{thm:no_productHall}. \green{} The base rank~$1$ is closed by the
same separating-scale arithmetic (\code{rank\_one\_coreCollision\_absurd}), with no external SNOL.
Strong induction over the rank then gives collision $\Rightarrow$ \code{Engine}, and a pigeonhole on
the finite signature (\code{coreCollision\_of\_infinite}) manufactures the collision from an infinite
injective family of starts.

\begin{theorem}[\code{product\_core\_engine\_of\_carrier}]\label{thm:product_core_engine_of_carrier}
Under the separating scale, $1\le r\le4$, and an infinite $S$ of ambient-legal \code{RankNode}s with
$\mathrm{node}$ injective on $S$, we obtain \code{Engine}.
\end{theorem}

\emph{Proof idea.} Assemble the extensionality, the bound, finiteness, the descent
(\cref{thm:core_step_proved}), the arithmetic base and the pigeonhole
(\code{coreCollision\_of\_infinite}). \green{} A \emph{conditional} theorem: everything between the
carrier hypothesis and the \code{Engine} is machine-checked. The hole is compressed into one input,
\code{CarrierData} --- an infinite injective family of ambient-legal starts of fixed rank.

\subsection{Factorisation and the last link}

The carrier hypothesis is supplied, node by node, by pure factorisation arithmetic.

\begin{theorem}[\code{mkNode\_of\_composite}]\label{thm:mkNode_of_composite}
Let $1\le A$, $1<N$, $N$ composite, $N\le 6X_A+1$, $6X_A+1<A^5$, and every prime divisor of $N$ greater
than $A$. Then for any sign there is a rank $r$ with $2\le r\le4$ and $X:\code{RankNode}\,r$ with
$\code{AmbientLegal}\,X_A\,X.\mathrm{factors}$ and $\prod(\mathrm{ofFn}\,X.\mathrm{factors})=N$.
\end{theorem}

\emph{Proof idea.} Take $L=\code{primeFactorsList}\,N$: compositeness gives $|L|\ge2$
(\code{composite\_rank\_ge\_two}); cleanliness makes every factor $>A$; and $(A+1)^{|L|}\le\prod L=N\le
6X_A+1<A^5<(A+1)^5$ forces $|L|\le4$ (\code{factor\_rank\_le\_four}). The witness $N_0=N$ certifies
\code{AmbientLegal}. \green{} The magic number $4$ is an arithmetic consequence of the sieve scale, not
a convention. Joining this node factory to the proven infinitude of the carrier:

\begin{theorem}[\code{cleanCenters\_infinite}]\label{thm:cleanCenters_infinite}
For every $A$, the set $\mathrm{CleanCenters}(A)=\{m\mid\mathrm{CleanZ}\,A\,m\}$ is infinite.
\end{theorem}

\emph{Proof idea.} By \code{Set.infinite\_of\_not\_bddAbove}: for any $N$, \code{carrier\_nonempty\_above}
produces a clean $m>N$ (a multiple of the old primorial). \green{} No density is needed.

\begin{theorem}[\code{engine\_of\_carrier\_and\_factorize}]\label{thm:engine_of_carrier_and_factorize}
Under the separating scale, an infinite carrier $C$, and a rank map \code{rankOf} and node map
\code{mkNode} injective and \code{AmbientLegal} on each rank, \code{Engine} holds.
\end{theorem}

\emph{Proof idea.} The composition of \code{factorizationData\_of\_carrier} (an infinite fibre of one
rank, chosen by the infinite pigeonhole \code{exists\_infinite\_fiber}) with
\cref{thm:product_core_engine_of_carrier}. \green{} The infinitude of the carrier, the fibre choice,
and the whole pump machine are proved in full. The single input is the factorisation maps.

\subsection{The single open node}

Here we must be exact, or a reduction gets passed off as a proof. Legality $\code{AmbientLegal}\,X_A$
requires the factors to divide a common top side $N_0\le6X_A+1$, i.e. $m\lesssim X_A$: a \emph{finite}
horizontal window. But \cref{thm:cleanCenters_infinite} gives infinitude \emph{vertically},
$m\to\infty$, inevitably above $X_A$, where the certificate --- and with it the rank bound $\le4$ ---
fails. The two proven ends do not splice at a fixed scale. What is required is a mechanism supplying
infinitely many legal nodes at a \emph{growing} scale, with $A$ and $X_A$ growing together with $m$.

\begin{conjecture}[\code{GlobalOldAbsorption}, the single open node]\label{conj:globaloldabsorption}
\red{} There exists an infinite family of clean centres $m$, coherently as $m\to\infty$, each with a
composite side admitting a legal factorisation of fixed rank $1\le r\le4$ into primes $>A(m)$ under a
scale $A(m),X_A(m)$ preserving the separating scale $6X_A(m)+1<P_{A(m)}$, with distinct centres
yielding distinct nodes.
\end{conjecture}

This is the last Step00 obligation. It is \emph{irreducible} --- in strength equal to the twin
conjecture itself --- and it cannot be closed from inside the system: closing it is exactly running the
engine whose impossibility is the foundation. The barrier is structural (the horizontal window cuts off
any fixed scale), not an artifact of one route: counting, EPMI descent, and the four-corner all
converge here, and every wrapper attempt was machine-audited back to it. Accordingly the machine
records exactly one \code{sorry}-tainted declaration, \code{twin\_prime\_conjecture}, closing through
\cref{thm:engine_of_carrier_and_factorize} onto this node.

\begin{remark}
We restate the honest verdict. The chain
\[
  \code{survivor\_of\_not\_covered}\to\code{twin\_prime\_conjecture\_of\_blocks}\to
  \code{twin\_primes\_of\_nopsl}\to\code{product\_core\_engine\_of\_carrier}\to
  \code{engine\_of\_carrier\_and\_factorize}
\]
is \green{} machine-checked and distribution-free from end to end. The twin conjecture is
\emph{not} proved. What is proved is the \emph{reduction} of it to the single node
\cref{conj:globaloldabsorption}, marked \red{}, which \cref{sec:firstcause} accepts as the first
cause \axiom{} from outside the system.
\end{remark}
